\section{Herramienta 1: Interpolación}

\begin{frame}{¿Qué es la interpolación?}
  \begin{block}{Definición intuitiva}
    Es una herramienta matemática que permite \textbf{estimar valores desconocidos a partir de valores conocidos}.
  \end{block}
  \vspace{0.5em}
  \begin{itemize}
    \item Si tienes datos en los minutos 5, 10 y 15, puedes estimar qué pasa en el minuto 7 o 12.
    \item Se construye una curva que pasa exactamente por los puntos medidos.
  \end{itemize}
\end{frame}

\begin{frame}{Ejemplo concreto en Dota 2}
  \begin{columns}
    \column{0.45\textwidth}
    \begin{center}
      \begin{tabular}{cc}
        \toprule
        Minuto & Oro \\
        \midrule
        5  & 1200 \\
        10 & 2800 \\
        15 & 4500 \\
        \bottomrule
      \end{tabular}
    \end{center}
    \column{0.55\textwidth}
    \begin{tikzpicture}
      \begin{axis}[
        width=6.5cm, height=5cm,
        xlabel={Minuto}, ylabel={Oro},
        xmin=0, xmax=18, ymin=0, ymax=5200,
        grid=major, grid style={gray!20},
        tick label style={font=\scriptsize},
        label style={font=\scriptsize},
      ]
        \addplot[only marks, mark=*, optcolor, mark size=2.5pt] coordinates {
          (5,1200) (10,2800) (15,4500)
        };
        \addplot[optcolor, thick, smooth] coordinates {
          (5,1200) (7,1850) (10,2800) (12,3550) (15,4500)
        };
        \addplot[only marks, mark=triangle*, realcolor, mark size=3pt] coordinates {
          (12,3550)
        };
        \node[font=\scriptsize, realcolor] at (axis cs:13,3300) {estimado};
      \end{axis}
    \end{tikzpicture}
  \end{columns}
  \vspace{0.3em}
  \begin{itemize}
    \item No medimos el minuto 12: la interpolación lo estima.
  \end{itemize}
\end{frame}

\begin{frame}{¿Por qué interpolación y no otra herramienta?}
  \begin{itemize}
    \item Los datos del juego son \textbf{discretos} (minuto a minuto).
    \item El tiempo del juego es \textbf{continuo} (fluye segundo a segundo).
    \item La interpolación \textbf{llena los huecos} entre los datos que sí tenemos.
  \end{itemize}
  \vspace{0.8em}
  \begin{block}{En una frase}
    La interpolación convierte datos sueltos en una función continua que podemos analizar.
  \end{block}
\end{frame}